\section{采样分配的完整解释}
\subsection{组指标与位置指标分别起什么作用}
联合策略存在两级决定。组指标先决定下一轮新增资源在左内部、右内部和界面之间如何分配；位置指标再决定每组内部选哪些候选点。组平均残差较大，并不意味着该组内所有位置都同样困难；某个位置残差很大，也不意味着把整轮预算全部给该组最合适。这两级结构需要分开记录，否则最终误差无法对应到具体采样行为。

本文组指标使用RMS，既考虑高残差幅值，也受高残差覆盖比例影响。少数极端点可能提高指标，而一个范围较广但中等残差的区域也可能具有较大RMS。将其改为最大值、分位数或平方均值都会改变分配规则。本文未把这些备选逐一优化，只固定当前定义进行检验；不把实验后想到的指标写成已经测试的方法。

界面指标把值误差与按最大系数归一化的通量误差组合。它具有简单、无额外训练的优点，但隐含一个数值优先级：对于高系数侧，较大的物理通量误差可能被缩小到与值误差相近的量级。实验因此同时报告归一化和物理通量，不能用前者很小替代后者满足应用精度要求。

\subsection{预算分配的代数形式}
设三组指标为$\eta_k$，最低点数为$b_k$，每个物理位置成本为$a_k$，本轮成本预算为$B_0$。先令剩余成本$R=B_0-\sum_ka_kb_k$。当$R\ge0$且指标不全近零时，实数新增目标为
\[
t_k=\frac{R(\eta_k+\varepsilon)}{a_k\sum_j(\eta_j+\varepsilon)}.
\]
先取整数$\lfloor t_k\rfloor$，再在剩余成本可容纳的组间补齐。若指标均接近零，改用相等组权重，避免除以微小总量。实际代码中的阈值为$10^{-14}$，属于数值保护而非误差容限。

成本为1的组使本实验整数预算总能填满；对任意成本向量并不具有这个保证。例如所有可选成本均为2而剩余成本为1时无法填满，必须改成不超过预算或重新定义允许成本。本文算法对本研究两种成本约定合法，不宣称解决一般整数规划问题。

最低配额会在指标非常不平衡时限制集中程度，避免某侧长期完全不新增点，但也会消耗可用于困难区域的预算。这一选择是覆盖与集中之间的工程折中。既有数据没有证明$(16,16,8)$是最优配额，也没有考察把它改成随轮次变化的规则；应按真实设置理解结果。

\subsection{抽样概率并不等于最终点密度}
探索比例0.2使可用池内每点在当前批次开始时具有非零选择权重。它不表示最终训练集中恰有20\%的点均匀分布，因为无放回过程、历史点、组配额和多轮更新共同决定最终集合。同样，最终点数多的区域不一定在每轮都具有最大的残差。

历史点保留可避免丢弃已经满足的局部约束，但也使早期选点持续影响后续目标。若采用替换策略，每轮更新后的训练测度和求解成本会不同。本文只实现增量路线，并在比较时把所有历史求解计入时间。不能把最终一次拟合时间当成完整自适应成本，也不能把与增量路线不同的已有采样方法视为逐公式复现。

候选池大小还限制能发现的结构。若窄峰附近没有足够候选点，残差排序无法选择池外位置；若池内存在相关或近重复位置，抽样无放回也不等于几何上有良好间距。均匀探索降低极端集中的风险，但不提供填充距离或连续覆盖保证。局部峰失败因此需要同时考虑候选分辨率与特征空间，不能只讨论界面预算。

\section{实验设计与结论强度}
\subsection{训练、候选、监测和最终测试}
训练点决定矩阵，候选点决定自适应更新，最终测试点只计算误差。三者用途不同。候选池虽在最初生成时与初始训练分开，其中部分位置后来会被选入训练；随着算法推进，候选指标已经受到选择历史影响，不能继续把它当作未接触训练的泛化测试。

D\_unused排除已经选用的候选位置后计算组指标，有助于诊断训练点混入带来的乐观性，但未用池仍经过历史选择筛选，故也不是新的独立样本。一个更严格的后续方案是在每轮另生成仅用于监测的池，且不从中挑点；若又根据监测误差调参数，则还需要保留未用于选择的最终测试。本文没有把这一后续设计冒充已经执行的实验。

现有测试网格与采样点的分离避免了最直接的训练误差自评，但探索性研究反复查看同一组测试结果仍可能形成研究者层面的选择偏差。因此本文明确区分主比较与后续诊断，不把后续扫描中表现较好的参数当成主实验事先确定的最优方案。真正的外推检验需要新模型或冻结规则后的新增运行。

\subsection{共享随机种子与统计单位}
同一设置内，A、B、C、D使用相同种子生成特征，这使配对差异较少受到完全不同空间的影响。不过每种策略的点集仍不同，种子共享只构成控制条件之一。跨系数组合重复相同种子时，相应特征空间高度相关，不能把三组五种子简单当作十五个独立空间。

本文以每个固定模型、预算和参数组合为一个设置，在其中报告五次重复及配对比值。跨设置合并的中位数只作为描述摘要，不附上假设独立重复后才成立的显著性结论。五个种子也不足以准确估计罕见失败概率；没有看到失败不等于已证明失败概率为零。

标准差和均值的用途不同。均值反映该五种子集合的平均误差，标准差反映离散程度，最大值揭示是否存在明显较差初始化。对于数值跨度较大的误差，少数较差运行会影响算术均值；本文保留逐次CSV，并同时使用配对比例中位数，避免只呈现一种摘要而失去分布信息。

\subsection{结果判据不随预期移动}
正确性判据包括制造数据一致、导数一致、FDM收敛与预算守恒。自适应效果判据则是合理基线、公平资源、独立误差和完整成本，不能要求策略必胜，也不能在其失败后悄悄换指标宣称成功。PDE残差稍低但真解误差更高时，应报告为取舍，不把原先的解精度目标事后改成仅优化残差。

数值状态ok只表示程序运行完成并返回有限指标。窄峰误差超过1说明当前配置没有达到高精度，即使最小二乘本身成功退出。本文把运行异常、低精度与未达到改进目标分别表述，因为三者的修正措施不同：异常需要程序排查，低精度需要模型或离散诊断，缺乏优势则可能说明当前方案本身不值得继续加复杂度。

\section{计算成本的可比较范围}
相同物理点数只控制一种资源，界面位置产生两行的事实使行数仍可能不同；相同行数又不控制几何覆盖。本文同时记录二者，并指出增量算法还有候选评估和反复求解开销。因此固定最终行预算的比较回答的是相近最终线性系统规模下的效果，不等价于严格相同总CPU时间。

在CPU稠密线性代数实现中，矩阵宽度增长通常比少量新增行更昂贵，奇异值诊断本身也有成本。本文计时包含这些诊断，是可复核的实现约定；若将来只用QR而不重复诊断，运行时间可能改变，但必须对全部方法一致调整。当前不能仅凭复杂度阶或某次耗时推断另一台机器的加速倍数。

进一步的精度目标比较应预先规定误差阈值、允许参数和搜索成本，报告达到阈值所需的完整时间。现有实验主要为固定配置与预算比较，尚未完成统一精度目标下的完整优化搜索。本文因此没有声称已经证明某方法在任意目标精度下更高效，也没有用不同算法各自最好的一次运行代替公平成本分析。
